\subsection{Depth and normality}

Let A be a Noetherian ring and p a prime ideal of A. We have \(\text{prof}(A_{\mathfrak{p}}) \leqslant \text{ht}(\mathfrak{p})\) (n° 4, cor. 2 of th. 2). If moreover A is reduced, the local ring \(A_{p}\) is reduced (II, § 2, n° 6, prop. 17), which implies:

\begin{enumerate}
\def\labelenumi{\alph{enumi})}
\item
  if ht(p) = 0, A \(_{p}\) is a field;
\item
  if \(\operatorname{ht}(\mathfrak{p}) \geqslant 1\) , one has \(\operatorname{prof}(A_{\mathfrak{p}}) \geqslant 1\) ( \(n^{\circ}\) 1, remark 3).
\end{enumerate}

Conversely:

PROPOSITION 15.--- Let A be a Noetherian ring satisfying the following two conditions:

\begin{enumerate}
\def\labelenumi{(\roman{enumi})}
\item
  for every minimal prime ideal \(\mathfrak{p}\) of A, the ring \(A_{\mathfrak{p}}\) is reduced;
\item
  for every prime ideal \(\mathfrak{p}\) of A of height \(\geqslant 1\) , one has \(\operatorname{prof}(A_{\mathfrak{p}}) \geqslant 1\) . Then A is reduced.
\end{enumerate}

Let \(\mathfrak{n}\) be the nilradical of A. For every minimal prime ideal \(\mathfrak{p}\) of A, we have by (i) \(\mathfrak{n}_{\mathfrak{p}} = 0\) , that is \(\mathfrak{p} \notin \operatorname{Supp}(\mathfrak{n})\) , and a fortiori \(\mathfrak{p} \notin \operatorname{Ass}_{\Lambda}(\mathfrak{n})\) . For every \(\mathfrak{p} \in \operatorname{Spec}(\mathbf{A})\) of height \(\geqslant 1\) , we have by (ii) \(\mathfrak{pA}_{\mathfrak{p}} \notin \operatorname{Ass}_{\mathrm{A}_{\mathfrak{p}}}(\mathrm{A}_{\mathfrak{p}})\) , and a fortiori \(\mathfrak{pA}_{\mathfrak{p}} \notin \operatorname{Ass}_{\mathrm{A}_{\mathfrak{p}}}(\mathfrak{n}_{\mathfrak{p}})\) , whence \(\mathfrak{p} \notin \operatorname{Ass}_{\mathrm{A}}(\mathfrak{n})\) (IV, § 1, no. 2, prop. 5). Thus the set \(\operatorname{Ass}_{\mathrm{A}}(\mathfrak{n})\) is empty, which implies that \(\mathfrak{n}\) is zero.

PROPOSITION 16.--- Let A be a Noetherian integrally closed ring, J an ideal of A of height ≥ 2, and M a reflexive finitely generated A-module. Then prof \(_{A}\) (J; M) ≥ 2.

Let us choose a finite-dimensional vector space V over the field of fractions of A and a lattice N in V isomorphic to M (VII, § 4, No. 2, Remark 1). The associated prime ideals of V/N, being of height 1 (loc. cit., Th. 2), do not contain J; by Remark 2 of No. 1, this implies \(\operatorname{prof}_{\mathrm{A}}(\mathrm{J};\mathrm{V}/\mathrm{N}) \geqslant 1\). On the other hand, the A-module V is divisible and torsion-free, hence injective (A, X, p. 17, Cor. 2 of Prop. 10), which implies \(\operatorname{prof}_{\mathrm{A}}(\mathrm{J};\mathrm{V}) = +\infty\). The inequality \(\operatorname{prof}_{\mathrm{A}}(\mathrm{J};\mathrm{N}) \geqslant 2\) then follows from Prop. 1 of No. 1.

COROLLARY.--- A Noetherian local ring that is integrally closed and of dimension ≥ 2 has depth ≥ 2.

Let A be a normal ring (no. 8) and let p be a prime ideal of A. Then:

\begin{enumerate}
\def\labelenumi{\alph{enumi})}
\item
  if ht(p) = 0, A \(_{p}\) is a field;
\item
  if \(\operatorname{ht}(\mathfrak{p}) = 1\) , \(A_{p}\) is a discrete valuation ring (VII, § 1, no. 7, prop. 11);
\item
  If \(\operatorname{ht}(\mathfrak{p}) \geqslant 2\), one has \(\operatorname{prof}(A_{\mathfrak{p}}) \geqslant 2\) (Corollary of Prop. 16).
\end{enumerate}

Conversely:

THEOREM 4 (Serre).--- Let A be a Noetherian ring satisfying the following conditions:

\begin{enumerate}
\def\labelenumi{(\roman{enumi})}
\item
  for every minimal prime ideal p of A, the ring A \(_{p}\) is reduced ;
\item
  for every prime ideal p of A of height 1, the ring A \(_{p}\) is integrally closed ;
\item
  for every prime ideal \(\mathfrak{p}\) of A of height \(\geqslant 2\) , one has \(\operatorname{prof}(A_{\mathfrak{p}}) \geqslant 2\) .
\end{enumerate}

Then A is normal.

The goal is to prove that the ring \(A_{\mathfrak{p}}\) is integrally closed for every \(\mathfrak{p} \in \operatorname{Spec}(A)\), which we shall do by induction on \(\operatorname{ht}(\mathfrak{p})\). For \(\operatorname{ht}(\mathfrak{p}) \leqslant 1\), this follows from hypotheses (i) and (ii). Suppose therefore that \(\operatorname{ht}(\mathfrak{p}) \geqslant 2\) and that \(A_{\mathfrak{q}}\) is integrally closed for every prime ideal \(\mathfrak{q}\) of A with \(\operatorname{ht}(\mathfrak{q}) < \operatorname{ht}(\mathfrak{p})\). By hypothesis (iii), we have \(\operatorname{prof}(A_{\mathfrak{p}}) \geqslant 2\). By the induction hypothesis and Corollary 3 of Theorem 3 of No. 9, applied to the ring \(A_{\mathfrak{p}}\) and to the closed subset \(\{\mathfrak{p}A_{\mathfrak{p}}\}\) of \(\operatorname{Spec}(A_{\mathfrak{p}})\), the ring \(A_{\mathfrak{p}}\) is an integral domain. Let K be its field of fractions, and let B be a subring of K containing \(A_{\mathfrak{p}}\) and finite over \(A_{\mathfrak{p}}\). It remains to prove that B is equal to \(A_{\mathfrak{p}}\). Let i denote the canonical injection from \(A_{\mathfrak{p}}\) into B. Since B is contained in K, it is an \(A_{\mathfrak{p}}\)-module without torsion, hence of depth \(\geqslant 1\). Moreover, for every prime ideal \(\mathfrak{q}\) of \(A_{\mathfrak{p}}\) distinct from \(\mathfrak{p}A_{\mathfrak{p}}\), the homomorphism \(i_{\mathfrak{q}}: (A_{\mathfrak{p}})_{\mathfrak{q}} \to B_{\mathfrak{q}}\) is bijective since \(A_{\mathfrak{q}}\) is integrally closed. By Lemma 4 of No. 9, applied to the closed subset \(F = \{\mathfrak{p}A_{\mathfrak{p}}\}\) of \(\operatorname{Spec}(A_{\mathfrak{p}})\), the homomorphism i is bijective, which completes the proof of the theorem.

Remark 1.--- A convenient form of Th. 4 is the following: let A be a Noetherian ring such that for every prime ideal p of A the ring A \(_{p}\) is integrally closed or of depth ≥ 2; then A is normal. Indeed, let us verify the hypotheses of Th. 4. Let p ∈ Spec(A). If ht(p) ≤ 1, then we have prof(A \(_{p}\) ) ≤ 1, hence A \(_{p}\) is integrally closed. If ht(p) ≥ 2, the ring A \(_{p}\) is either of depth ≥ 2, or integrally closed, which again implies depth(A \(_{p}\) ) ≥ 2 (no. 1, cor. 2 of prop. 1), whence the desired conclusion. One verifies similarly the following statement: let A be a Noetherian ring such that for every prime ideal p of A, the ring A \(_{p}\) is reduced or of depth ≥ 1; then A is reduced.

COROLLARY 1.--- Let \(A\) be a Noetherian ring, F a closed subset of \(\operatorname{Spec}(A)\), U the complementary open set. We assume that \(\operatorname{prof}_{\mathrm{F}}(\mathrm{A})\) is \(\geqslant 2\) (resp. \(\geqslant 1\) and that, for every \(\mathfrak{p} \in \mathrm{U}\), the ring \(\mathrm{A}_{\mathfrak{p}}\) is integrally closed (resp. reduced). Then A is normal (resp. reduced).

For every \(\mathfrak{p} \in \mathrm{F}\) , one has \(\operatorname{prof}(\mathrm{A}_{\mathfrak{p}}) \geqslant \operatorname{prof}_{\mathrm{F}}(\mathrm{A})\) ( \(n^{\circ}\) 5, prop. 8); it therefore suffices to apply the preceding remark.

COROLLARY 2.-- Let ρ : A → B be a homomorphism of Noetherian rings making B a flat A-module.

\begin{enumerate}
\def\labelenumi{\alph{enumi})}
\item
  If B is normal and faithfully flat over A, then A is normal.
\item
  Suppose that A is normal and that the ring \(\kappa(\mathfrak{p})\otimes_{\mathrm{A}}\mathrm{B}\) is normal when \(\mathfrak{p}\) is a minimal prime ideal of A, and reduced when \(\mathfrak{p}\) is a prime ideal of height 1. Then the ring B is normal.
\item
  Suppose B is normal and faithfully flat over A. Then \(\rho\) is injective (I, § 3, no. 5, prop. 9) and B is reduced, hence A is reduced. Let S be the set of elements of A that are not zero divisors. Since B is flat over A, \(\rho(S)\) is formed of non-zero-divisors in B, so that B is integrally closed in \(\rho(S)^{-1}B\) . Let \(x\) be an element of \(S^{-1}A\) integral over A; then the element \(x \otimes 1_B\) of the ring \(S^{-1}A \otimes_A B\) (which identifies with \(\rho(S)^{-1}B\) ) is integral over B, hence belongs to B. Since B is faithfully flat over A, we have \(x \in A\) (I, § 3, no. 5, prop. 10, (ii)), and A is normal.
\item
  Under the hypotheses of b), it suffices by Remark 1 to prove that for every prime ideal \(\mathfrak{q}\) of B, the local ring \(B_{\mathfrak{q}}\) is normal or has depth \(\geqslant 2\). Let \(\mathfrak{p}\) denote the prime ideal \(\rho^{-1}(\mathfrak{q})\) of A; the local homomorphism \(A_{\mathfrak{p}} \to B_{\mathfrak{q}}\) induced by \(\rho\) makes \(B_{\mathfrak{q}}\) a faithfully flat \(A_{\mathfrak{p}}\)-module (I, § 3, No. 5, Prop. 9). We distinguish four cases:
\end{enumerate}

\begin{enumerate}
\def\labelenumi{\arabic{enumi})}
\item
  If \(\operatorname{ht}(\mathfrak{p}) = 0\), then \(A_{\mathfrak{p}}\) is a field, equal to \(\kappa(\mathfrak{p})\); the ring \(A_{\mathfrak{p}} \otimes_{A} B\) is normal by hypothesis. The same holds for \(B_{\mathfrak{q}}\), which is a ring of fractions thereof.
\item
  If \(\operatorname{ht}(\mathfrak{p}) = 1\) and \(\operatorname{ht}(\mathfrak{q}) \leqslant 1\), then \(A_{\mathfrak{p}}\) is a discrete valuation ring; let \(\pi\) be a uniformizer of \(A_{\mathfrak{p}}\). Since \(B_{\mathfrak{q}}\) is faithfully flat over \(A_{\mathfrak{p}}\), the element \(\pi 1_{B_{\mathfrak{q}}}\) of \(B_{\mathfrak{q}}\) is not a zero divisor, so that the local ring \(B_{\mathfrak{q}}/\pi B_{\mathfrak{q}}\) has dimension 0 (VIII, § 3, No. 1, Cor. 2 of Prop. 1). It is reduced, since it is a ring of fractions of the reduced ring \(\kappa(\mathfrak{p}) \otimes_{A} B\), and consequently it is a field. Therefore \(B_{\mathfrak{q}}\) is a discrete valuation ring (VI, § 1, No. 4, Prop. 2), hence integrally closed.
\item
  \(\operatorname{ht}(\mathfrak{p}) = 1\) and \(\operatorname{ht}(\mathfrak{q}) \geqslant 2\) . Then, by the relation
\end{enumerate}

\[
\dim (B _ {\mathfrak {q}}) = \dim (A _ {\mathfrak {p}}) + \dim (\kappa (\mathfrak {p}) \otimes_ {A} B _ {\mathfrak {q}})
\]

(VIII, § 3, no. 4, cor. 1 of prop. 7), the ring \(\kappa(\mathfrak{p})\otimes_{\mathrm{A}}\mathrm{B}_{\mathfrak{q}}\) has dimension \(\geqslant 1\) . It is reduced by hypothesis, hence of depth \(\geqslant 1\) (no. 1, remark 3). We then have (no. 6, cor. of prop. 11)

\[
\operatorname{prof} \left(\mathrm{B} _ {\mathfrak {q}}\right) = \operatorname{prof} \left(\Lambda_ {\mathfrak {p}}\right) + \operatorname{prof} \left(\kappa (\mathfrak {p}) \otimes_ {\mathrm{A}} \mathrm{B} _ {\mathfrak {q}}\right) \geqslant 1 + 1 = 2.
\]

\begin{enumerate}
\def\labelenumi{\arabic{enumi})}
\setcounter{enumi}{3}
\tightlist
\item
  \(\operatorname{ht}(\mathfrak{p}) \geqslant 2\) . Since \(A_{p}\) is of depth \(\geqslant 2\) (corollary of Prop. 16), the same holds for \(B_{q}\) by loc. cit.
\end{enumerate}

COROLLARY 3.--- Let \(\rho: A \to B\) be a homomorphism of Noetherian rings. Suppose that \(B\) is a flat A-module, that \(A\) is normal, and that \(\kappa(\mathfrak{p}) \otimes_A B\) is normal for all \(\mathfrak{p} \in \operatorname{Spec}(A)\). Then \(B\) is normal.

